\chapter{Application Routing, Authentication, and Authorization}

\section{Route / Blueprint Architecture}

The application registers ten Flask blueprints in \code{create\_app()}
(\file{app/\_\_init\_\_.py}), plus one inline route (\code{/healthz})
registered directly on the app for container health checks. The table
below is a complete map of every route in the codebase, extracted
directly from the blueprint files.

\begin{longtable}{@{}p{1.5cm}p{2.9cm}p{1.0cm}p{1.5cm}p{4.2cm}p{3.6cm}@{}}
\toprule
\textbf{Blueprint} & \textbf{Route} & \textbf{Method} & \textbf{Role} &
\textbf{Key service call(s)} & \textbf{Template} \\
\midrule
\endhead
main & \route{/} & GET & any & --- (redirect) & redirects to \longcode{dashboard.index} \\
auth & \route{/login} & GET/POST & public & \longcode{auth.service.authenticate} & \file{auth/login.html} \\
auth & \route{/logout} & POST & any & \longcode{flask\_login.logout\_user} & redirect \\
dashboard & \route{/dashboard} & GET & any & \longcode{reports.*} (many) & \file{dashboard/admin.html} or \file{dashboard/employee.html} \\
dashboard & \route{/reports/overtime} & GET & admin, manager & \longcode{reports.overtime\_summary} & \file{reports/overtime.html} \\
dashboard & \route{/reports/hours-trend} & GET & admin, manager & \longcode{reports.hours\_trend} & \file{reports/hours\_trend.html} \\
attendance & \route{/attendance} & GET & any & \longcode{reports.attendance\_entries\_with\_context}, \longcode{reports.current\_attendance\_status} & \file{attendance/list.html} \\
attendance & \route{/attendance/clock-in} & POST & any & \longcode{attendance.clock\_in} & redirect \\
attendance & \route{/attendance/\allowbreak<id>/clock-out} & POST & any & \longcode{attendance.clock\_out} & redirect \\
attendance & \route{/attendance/\allowbreak<id>/correct} & POST & admin, manager & \longcode{attendance.correct\_entry} & redirect \\
schedule & \route{/schedule} & GET & any & \longcode{scheduling.list\_shifts} & \file{schedule/list.html} \\
schedule & \route{/schedule} & POST & admin, manager & \longcode{scheduling.create\_shift} & redirect / re-render \\
schedule & \route{/schedule/<id>} & POST & admin, manager & \longcode{scheduling.update\_shift} & redirect \\
schedule & \route{/schedule/\allowbreak<id>/assign} & POST & admin, manager & \longcode{scheduling.assign\_employee} & redirect \\
schedule & \route{/schedule/\allowbreak<id>/publish} & POST & admin, manager & \longcode{scheduling.publish\_shift} & redirect \\
schedule & \route{/schedule/\allowbreak<id>/cancel} & POST & admin, manager & \longcode{scheduling.cancel\_shift} & redirect \\
leave & \route{/leave} & GET & any & \longcode{leave.list\_leave\_requests} & \file{leave/list.html} \\
leave & \route{/leave} & POST & any & \longcode{leave.request\_leave} & redirect \\
leave & \route{/leave/\allowbreak<id>/approve} & POST & admin, manager & \longcode{leave.approve\_leave} & redirect \\
leave & \route{/leave/<id>/reject} & POST & admin, manager & \longcode{leave.reject\_leave} & redirect \\
leave & \route{/leave/<id>/cancel} & POST & any & \longcode{leave.cancel\_leave} & redirect \\
employees & \route{/employees} & GET & admin, manager & \longcode{employees.list\_employees} & \file{employees/list.html} \\
employees & \route{/employees} & POST & admin & \longcode{employees.create\_employee} & redirect / re-render \\
employees & \route{/employees/<id>} & GET & any & \longcode{employees.get\_employee}, \longcode{employees.get\_linked\_user} & \file{employees/detail.html} \\
employees & \route{/employees/<id>} & POST & admin, manager & \longcode{employees.update\_employee} & redirect \\
employees & \route{/employees/\allowbreak<id>/terminate} & POST & admin & \longcode{employees.terminate\_employee} & redirect \\
employees & \route{/employees/\allowbreak<id>/pay-rate} & GET & admin & \longcode{pay\_rates.list\_pay\_rate\_history} & \file{employees/pay\_rate.html} \\
employees & \route{/employees/\allowbreak<id>/pay-rate} & POST & admin & \longcode{pay\_rates.set\_pay\_rate} & redirect \\
departments & \route{/departments} & GET & admin, manager & \longcode{departments.list\_departments} & \file{departments/list.html} \\
departments & \route{/departments} & POST & admin & \longcode{departments.create\_department} & redirect / re-render \\
departments & \route{/departments/<id>} & POST & admin & \longcode{departments.update\_department} & redirect \\
departments & \route{/departments/\allowbreak<id>/deactivate} & POST & admin & \longcode{departments.deactivate\_department} & redirect \\
labor\_cost & \route{/labor-cost} & GET & admin, manager & \longcode{labor\_cost.department\_cost\_summary} & \file{labor\_cost/index.html} \\
labor\_cost & \route{/labor-cost/\allowbreak employees/<id>} & GET & admin & \longcode{labor\_cost.range\_cost\_for\_employee} & \texttt{\small\seqsplit{labor\_cost/employee\_detail.html}} \\
audit & \route{/audit-log} & GET & admin & \longcode{audit.list\_entries} & \file{audit/list.html} \\
audit & \route{/recent-activity} & GET & admin & \longcode{reports.recent\_activity} & \file{audit/recent\_activity.html} \\
(app-level) & \route{/healthz} & GET & public & --- & JSON \longcode{\{"status": "ok"\}} \\
\bottomrule
\end{longtable}

\noindent\textit{"Role" above is the minimum \code{role\_required(...)}
on the route; ``any'' means \code{login\_required} only (every
authenticated role may reach the route, though what it returns still
differs by role --- e.g.\ \route{/attendance} shows only the caller's
own entries for an employee).}

\section{Request Lifecycle}

Every page load in the application follows the same path. This is the
concrete version of the generic layered architecture from Chapter~1,
now including where \code{AccessScope} is built and where the
Jinja/CSS layer fits in.

\begin{center}
\begin{tikzpicture}[node distance=0.35cm]
  \node[flownode] (n1) {Browser sends an HTTP request (session cookie attached)};
  \node[flownode, below=of n1] (n2) {Flask routes the request to a blueprint view function};
  \node[flownode, below=of n2] (n3) {Flask-Login resolves \code{current\_user} from the session (\code{load\_user})};
  \node[flownode, below=of n3] (n4) {\code{login\_required} / \code{role\_required} decorator checks};
  \node[flownode, below=of n4] (n5) {\code{build\_scope\_for\_user(current\_user)} builds an \code{AccessScope}};
  \node[flownode, below=of n5] (n6) {Route calls one or more service functions, passing the scope};
  \node[flownode, below=of n6] (n7) {Service function re-checks authorization, then queries via SQLAlchemy};
  \node[flownode, below=of n7] (n8) {PostgreSQL executes the query, enforcing its own constraints};
  \node[flownode, below=of n8] (n9) {Service returns plain data (models / dataclasses) to the route};
  \node[flownode, below=of n9] (n10) {Route calls \code{render\_template(...)} with that data};
  \node[flownode, below=of n10] (n11) {Jinja renders HTML, pulling in CSS/JS from \file{static/}};
  \node[flownode, below=of n11] (n12) {Flask sends the HTML response, with security headers attached};

  \foreach \a/\b in {n1/n2,n2/n3,n3/n4,n4/n5,n5/n6,n6/n7,n7/n8,n8/n9,n9/n10,n10/n11,n11/n12}
    \draw[flowarrow] (\a) -- (\b);
\end{tikzpicture}
\end{center}

The security headers in the last step (\code{X-Content-Type-Options},
\code{X-Frame-Options}, \code{Referrer-Policy}, a strict
\code{Content-Security-Policy}, and conditionally \code{Strict-\allowbreak
Transport-Security}) are attached by an \code{@app.after\_request} hook
in \file{app/\_\_init\_\_.py} to \emph{every} response, not just
authenticated pages. See Chapter~6, \S6.1, for why the CSP can afford to
be this strict (no external scripts, no inline styles, no CDN fonts).

\section{Authentication}

\subsection*{Login}

\code{POST /login} (\file{app/routes/auth.py}) delegates credential
checking entirely to \longcode{app.auth.service.authenticate}. That function:

\begin{enumerate}[nosep]
  \item Looks up the user by email. If none exists, the account is
        inactive, or the account is currently locked, it still runs a
        \emph{real} Argon2 verification against a fixed dummy hash
        (\code{\_burn\_dummy\_verification\_time}) before returning
        failure --- this keeps response timing statistically
        indistinguishable across every failure reason, closing a
        user-enumeration timing side-channel.
  \item On a genuine password mismatch, increments
        \code{failed\_login\_count} and, once it reaches
        \code{MAX\_FAILED\_LOGIN\_ATTEMPTS} (5), sets
        \code{locked\_until} to 15 minutes in the future
        (\code{LOCKOUT\_DURATION}).
  \item On success, resets the failure counter, records
        \code{last\_login\_at}, and stages an audit log entry.
  \item Every branch returns the identical, generic
        \code{"Invalid email or password."} message --- the UI never
        reveals \emph{which} of "no such account" / "wrong password" /
        "account locked" applies.
\end{enumerate}

On success, the route (\file{app/routes/auth.py}) calls
\code{session.clear()} \emph{before} \code{login\_user(...)} --- this
discards any pre-authentication session state (including the
already-validated CSRF token) so nothing from before login carries into
the authenticated session, a defense against session fixation. It then
sets \code{session.permanent = True}, so the cookie expires
\code{PERMANENT\_SESSION\_LIFETIME} (12 hours) after login regardless of
activity, and redirects to a same-site-validated \code{next} target
(\code{app.auth.redirects.get\_safe\_redirect\_target} --- see \S6.1
for the open-redirect defense).

\subsection*{Password Hashing}

Passwords are hashed with Argon2 (\code{argon2.PasswordHasher}'s
defaults) in \file{app/auth/passwords.py}. The plaintext password is
never stored, logged, or included in an audit entry anywhere in the
codebase (verified: \code{audit\_service.record}'s \code{changes} dict
is built explicitly, field by field, at every call site --- there is no
code path that could pass a whole form or a raw password through).

\subsection*{Sessions}

Flask-Login's signed session cookie is the only session mechanism ---
there is no separate server-side session table. This means the
\code{load\_user} hook in \code{app/models/user.py} (invoked on
\emph{every} request for an authenticated user) is the actual
enforcement point for "has this account been deactivated or locked out
since the cookie was issued?" --- it returns \code{None} (forcing
re-authentication) if \code{is\_active} is false or \code{locked\_until}
is still in the future, even for an otherwise validly-signed cookie.

\subsection*{Account Lockout}

Implemented exactly as described under Login above:
\code{failed\_login\_count} (a small integer, \code{CHECK ... >= 0}) and
\code{locked\_until} (nullable timestamptz) on the \code{users} table.
A wrong-password attempt against an already-\emph{expired} lock resets
the counter to 1 (treating it as the first attempt of a fresh window)
rather than incrementing a stale count back over the threshold --- see
the inline comment in \code{authenticate} for why the naive version of
this logic could re-lock an account indefinitely from one further wrong
guess per lockout period.

\subsection*{Unauthorized Requests}

An unauthenticated request to any \code{login\_required} route is
redirected to \code{/login} by Flask-Login (\code{login\_manager.login\_view
= "auth.login"}), with a flashed message ("Please log in to access this
page."). A request from an authenticated user lacking the required role
receives a plain \textbf{403} (\code{abort(403)} inside
\code{role\_required}), rendered by \file{templates/errors/403.html}.

\section{Authorization}

This is the most important subsystem in the codebase to understand
correctly, because a mistake here is a data-leak or privilege-escalation
bug, not a cosmetic one.

\subsection*{AccessScope}

Every authorization decision in the application is built from one
frozen dataclass, \code{app.auth.scope.AccessScope}:

\begin{lstlisting}[language=Python]
@dataclass(frozen=True)
class AccessScope:
    user_id: int
    organization_id: int
    role: str                       # 'admin' | 'manager' | 'employee'
    department_ids: frozenset[int]  # populated only for 'manager'
    employee_id: int | None         # populated for 'employee' (and any
                                     # admin/manager who also has a
                                     # linked employee record)
\end{lstlisting}

\code{build\_scope\_for\_user(current\_user)} constructs this once per
request, at the top of every route, by querying
\code{department\_managers} for a manager's assigned department ids.
Every route in the table in \S2.1 calls this before doing anything
else.

\subsection*{Defense in Depth}

Authorization is checked at up to three independent layers, and the
codebase's own module docstrings repeatedly call this out as
deliberate:

\begin{center}
\begin{tikzpicture}[node distance=0.4cm]
  \node[flownode] (a) {Route decorator (\code{login\_required} / \code{role\_required}) --- coarse, role-only};
  \node[flownode, below=of a] (b) {Service function re-checks role \emph{and} organizational/department scope};
  \node[flownode, below=of b] (c) {Database query is itself filtered by \code{organization\_id} (and, for a manager, \code{department\_id})};
  \node[flownode, below=of c] (d) {Composite foreign keys make a cross-tenant row physically unrepresentable};
  \node[flownode, below=of d, fill=white, draw=inkblue] (e) {\textbf{Authorized data reaches the template}};

  \foreach \a/\b in {a/b,b/c,c/d,d/e}
    \draw[flowarrow] (\a) -- (\b);
\end{tikzpicture}
\end{center}

\textbf{Why not rely on hiding UI elements?} A hidden button is a
client-side courtesy, not a security control --- a browser's developer
tools or a hand-crafted \code{curl} request can submit any form
regardless of what the rendered HTML shows. Every service function in
this codebase enforces its own rule regardless of what the calling
route or template does, so an attacker who reaches a route directly
(bypassing the UI entirely) gets exactly the same rejection a normal
user browsing the hidden feature would have. \file{tests/routes/test\_idor\_sweep.py}
exists specifically to exercise this: it drives routes directly with an
out-of-scope id and asserts a 404, not just checks that a link is
absent from a page.

\subsection*{The IDOR Defense: \code{get\_scoped\_or\_404}}

\code{app.auth.scope.get\_scoped\_or\_404(model, obj\_id, scope)} is the
standard way to fetch a single tenant-owned row. It always filters by
\code{organization\_id}, and additionally filters by
\code{department\_id} when the model has that column and the caller is
a manager. Critically, an out-of-scope row returns \textbf{404, not
403} --- so a caller cannot distinguish "this id doesn't exist" from
"this id exists but belongs to someone else," which would otherwise
leak the existence of other organizations'/departments' data through
the response code alone. Several models without a \code{department\_id}
column of their own (\code{AttendanceEntry}, \code{LeaveRequest})
implement the same 404-not-403 pattern by hand, via a join to the
owning employee (e.g.\ \code{attendance.\_get\_entry\_for\_scope}).

\subsection*{Organization Scoping}

Every tenant-owned table carries an \code{organization\_id} column, and
every service query filters on it. This is reinforced at the database
level: every child table's foreign key to a parent (e.g.\ an employee's
department, a shift's employee) is a \emph{composite} foreign key on
\code{(parent\_id, organization\_id)}, referencing a matching composite
unique constraint on the parent --- so it is not merely
application-code discipline that prevents, say, an employee row from
pointing at a department in a different organization; the schema itself
makes it impossible.

\subsection*{Manager Department Scoping}

A manager's \code{department\_ids} (Chapter~2's \code{AccessScope})
comes from the \code{department\_managers} join table. Every
manager-reachable service function filters its query by
\code{department\_id IN scope.department\_ids} (directly, or via a join
to the owning employee for tables with no department column). Because
this filter is applied \emph{before} any caller-supplied filter
parameter (e.g.\ \route{?employee\_id=} or \route{?department\_id=} on
several list views), a manager passing an out-of-scope id gets an
\emph{empty intersection}, never that other department's data --- this
composition order is exercised directly by tests such as
\code{test\_manager\_filtering\_by\_an\_unmanaged\_employee\_sees\_no\_data}.

\subsection*{Admin Privileges}

An admin has no \code{department\_ids} restriction at all --- every
manager-or-admin-reachable service function's manager-scoping branch is
skipped entirely for \code{role == "admin"}. Admin is still confined to
their own \code{organization\_id} like every other role (there is no
concept of a cross-organization "super-admin" anywhere in the schema or
code).

\subsection*{Manager Privileges}

Everything an admin can do within their managed departments, with two
explicit, hard-coded exceptions the codebase calls "rule A4": a manager
may see a department's aggregate labor cost \emph{total}
(\code{labor\_cost.department\_cost\_summary}) and an employee's
overtime \emph{hours} (\code{reports.overtime\_summary}, which reads
only \code{LineItem.hours}/\code{.category}, never \code{.rate} or
\code{.cost}) --- but never an individual employee's hourly pay rate or
per-employee cost breakdown. Those two views
(\route{/employees/<id>/pay-rate} and
\route{/labor-cost/employees/<id>}) are hard-coded
\code{role\_required("admin")}, not \code{role\_required("admin",
"manager")} like most of the rest of the application.

\subsection*{Employee Privileges}

Self-service only: their own schedule (published shifts only), their
own attendance (clock in/out, view history), their own leave requests.
An employee-role \code{AccessScope} has no \code{department\_ids} at
all (an empty frozenset) --- every service function's employee branch
checks identity (\code{scope.employee\_id == target\_id}) rather than
department membership.

\section{Role Matrix}

Verified directly against the \code{role\_required(...)} decorators and
each service function's own internal role check --- not guessed from
route names.

\begin{longtable}{@{}p{6.4cm}p{2.6cm}p{2.9cm}p{2.9cm}@{}}
\toprule
\textbf{Feature} & \textbf{ADMIN} & \textbf{MANAGER} & \textbf{EMPLOYEE} \\
\midrule
\endhead
View own dashboard & Org-wide overview & Managed-department overview & Personal dashboard \\
View / create / edit departments & Full & View only & No access \\
Deactivate a department & Yes & No & No \\
View employee list & Org-wide & Managed departments & No access (route is admin/manager only) \\
View one employee's detail page & Any employee & Managed-department employees & Own record only \\
Create / edit an employee & Yes & Edit only, managed departments & No \\
Terminate an employee & Yes & No & No \\
View / set an employee's pay rate & Yes & No & No \\
View a department's labor cost total & Yes & Managed departments & No \\
View an employee's labor cost breakdown & Yes & No & No \\
View / create shifts & Org-wide & Managed departments & View own published shifts only \\
Assign / publish / cancel a shift & Yes & Managed departments & No \\
Clock in / out (self) & Yes & Yes & Yes \\
Clock in/out on another's behalf & Yes & Managed departments & No \\
Correct an attendance entry & Yes & Managed departments & No (not even their own) \\
Request leave (self) & Yes & Yes & Yes \\
Request leave on another's behalf & Yes & Managed departments & No \\
Approve / reject leave & Yes & Managed departments & No \\
Cancel a leave request & Any & Managed departments & Own, while still pending \\
View overtime / hours-trend reports & Yes & Managed departments & No \\
View the audit log / recent activity & Yes & No & No \\
\bottomrule
\end{longtable}
